\def\ChapterCode{SHF}
\chapter
[\CO{[SHF]} Sehen\,--\,Hören\,--\,Fühlen]
{Sehen\,--\,Hören\,--\,Fühlen:\\Basis-Patienteneinschätzung und Einführung in die \enquote{Primary Survey}}
\label{chp:BET}

\ChapterIntro
{%
Am Ende des Kapitels soll der Anwender in der Lage sein,
den einen Patienten und seine Umgebung zu beschreiben, sowie
Probleme anhand seiner \sout{fünf} vier Sinne zu erkennen.
Es legt die Grundlage für das Kapitel Patientenmanagement und den
Einschätzungsblock, mittels welchem Pateinten im
professionellen Setting erstuntersucht werden.

}{%
\begin{description}
  \item[Maintainer:] Sebastian Gabriel
  \item[Co-Maintainer:] Josef Emhofer, Roman Koch
  \item[Autoren:] Diverse
  \item[Reviewer:] Standard-Reviewprozess
  \item[Version:] {\DocVersion} ({\DocVersionInfo})
  \item[SHA1:] {\DocGitCommit}
\end{description}

{\TxTeilkapitelAASS}
}

Grundregeln:

\begin{enumerate}
  \item Erst Ansehen,
  \item dann mit freien Ohren hören,
  \item dann evtl. riechen,
  \item und zum Schluss Tasten.
  \item \Z{Der Geschmackssinn bleibt dem Privatgebrauch
  vorbehalten.}
\end{enumerate}


\section{Umgebung und Situation}

Beobachten:

\begin{itemize}
  \item Lokalität: Freie Fläche, Verkehrsfläche, Werkstätte,
  Einfamilienhaus, Wohnung, Verkehrsmittel, Ordination
  \item Zufahrt
  \item Zeit: Tageszeit, Jahreszeit
  \item Witterung: Temperatur, Regen, Schnee, \ldots
  \item Anwesende Personen: Angehörige, Zeugen, Täter (!),
  andere Einsatzkräfte
  \item Sonstige Auffälligkeiten
\end{itemize}

Daraus kann dann abgeleitet werden:

\begin{itemize}
  \item Gefahren für sich
  
  Verkehr, Gefahrenzonen, Gefahrengüter, Haustiere, Gas, \ldots
  \item Gefahren für den Patienten
  
  Wie oben, zusätzlich z.\,B.: Unterkühlung, Überhitzung
  \item Erste Hinweise auf den Vorfall, Trauma bzw.
  Unfallmechanismus
  \item Anzahl der Patienten
\end{itemize}

\section{Der Patient}

\subsection{Patienten beobachten}

Das Gegenüber wird genau beobachtet. Zuerst betrachtet
man den ganzen Patienten und achtet besonders auf die
Körperhaltung:
\begin{itemize}
  \item Liegt der Patient? Sitzt er? Steht er?
  \item Ist seine Haltung locker oder angespannt?
  \item Nimmt er eine Schonhaltung ein? Vermeidet er
  bestimmte Bewegungen?
  \item Ist an den Körperteilen eine abnorme Stellung
  auffällig?
  \item Fallen sonst Verletzungen auf?
\end{itemize}

Im nächsten Schritt achtet man auf die Haut und das Gesicht:

\begin{itemize}
  \item Wie ist die Hautfarbe des Patienten? Gerötet Blass?
  Bläulich?
  \item Schwitzt der Patient?
  \item Nimmt der Patient Blickkontakt auf, oder nimmt er
  mich nicht wahr?
  \item Wie sieht sein Gesichtsausdruck aus?
\end{itemize}

Nachdem wir den Patienten ausführlich betrachtet haben,
nehmen wir Kontakt mit ihm auf. Wir stellen uns vor:

\begin{itemize}
  \item \Speech{\enquote{Guten Tag! Mein Name ist
  \Lang{\textit{Name}} von der/vom
  \Lang{\textit{Organisation}}.}}
\end{itemize}

Dabei können wir beobachten, ob er bewusstseinsklar (wach)
ist oder nicht. Ist dies der Fall, fragen wir ihn nach
seinen Beschwerden:

\begin{itemize}
  \item \Speech{\enquote{Was können wir für Sie tun?}}
\end{itemize}

Sofern der Patient weiß was sein Problem ist, wird er es nun
hoffentlich erzählen. 

\paragraph{Zwischenstand}

Wir haben eine neue Umgebung kennengelernt und beurteilt,
und haben nun einen Patienten vor uns, von dem wir
wissen, dass er bewusstseinsklar ist. Er hat uns auch seine Hauptbeschwerden genannt.

\subsection{Atemwege -- Atmung -- Kreislauf}

Als nächstes werden die grundlegenden Lebensfunktionen
beobachtet. 

\paragraph{Atemwege} 
Die Atemwege stehen als erste auf der
Liste: Sie können grundsätzlich frei oder verlegt sein. 


Bei Verdacht werden zuerst die \E{Mundhöhle}, die \E{Nase} und der
\E{Hals} \EE{inspiziert}, ob Fremdkörper oder Flüssigkeiten (Blut,
Erbrochenes) vorhanden sind.
Liegt ein Verdacht auf eine
Verletzung vor, werden Nase und Hals auch \E{abgetastet}.

Weiters werden die \EE{Atemgeräusche} beurteilt: 
Speziell \E{pfeifende}, \E{gurgelnde} oder \E{schnarchende}
Atemgeräusche können ein Hinweis auf eine Verlegung sein.
Oft gibt auch die \EE{Gestik} und Haltung des Patienten
Hinweise. 

\paragraph{Atmung} 
Sind die Atemwege frei, kann die Atmung
betrachtet werden. Dabei wird der \EE{Brustkorb betrachtet}, wie er sich hebt
und senkt. Man versucht ungefähr einzuschätzen \EE{wie oft}
der Patient atmet, und \EE{wie tief} diese Atmung ist.

Daneben gibt es noch andere Zeichen, die für die Atmung
interessant sind, zum Beispiel bedeutet eine bläuliche
Verfärbung. dass ein Sauerstoffmangel besteht. Die
\EE{Blaufärbung} sieht man besonders an den Fingern und den
Lippen gut.


\paragraph{Kreislauf}

Wir sehen uns nochmal die \EE{Hautfärbung} bewusst an: Ist
der Patient rosig oder eher bleich. Dann berühren wir ihn und
fühlen, ob sich die \EE{Haut} kühl \EE{anfühlt}. 

Als nächstes wird ein \EE{Rekap-Test} durchgeführt: Dabei
wird ein leichter Druck auf das Nagelende ausgeübt bis
das Nagelbett weiß ist. Dann lässt man los. Innerhalb von
\VonBis{1}{2} Sekunden sollte das Nagelbett wieder rosarot
sein. 



Dann wird der Radialis-Puls gefühlt un beurteilt, wie gut er
tastbar und wie schnell er ungefähr ist.




\paragraph{Dann gibt es noch \ldots}
einige Sachen die wir beurteilen können. 

Zum Beispiel die \EE{Pupillen}: Sind sie gleich weit? Rund?
Sind sie eng, weit oder mittelweit?

Weiters können wir noch fetsstellen, ob der Patient weiß wer
er ist, wo er ist, was gerade passiert und welche Zeit wir
haben:

\begin{itemize}
  \item \Speech{\enquote{Wie heissen Sie denn?}}
  \item \Speech{\enquote{Wissen Sie was gerade passiert ist?}}
  \item \Speech{\enquote{Können Sie mir sagen wo Sie gerade
  sind?}}
  \item \Speech{\enquote{Wissen Sie wie spät es ist?}}
  
  oder: \Speech{\enquote{Welche Tageszeit hebn wir?}}
  
  oder: \Speech{\enquote{Welche Jahreszeit haben wir?}}
  
  oder: \Speech{\enquote{Welches Jahr haben wir?}}
  
  Für Britische Staatsbürger: \Speech{\enquote{Who is the
  current Regent?}}
\end{itemize}


\section{Übersicht}


\begin{WidePar}
\FormatTable
\begin{tabularx}{\linewidth}[H]{p{2cm}XX}
% \toprule\addlinespace
&
\TabHeading{Beurteilung / Untersuchung}
&
\TabHeading{{\eye} Erkenne}
\\\addlinespace\toprule\addlinespace
\noindent\TabHeading{Szene}
&
\begin{minipage}[t]{\linewidth}
\begin{itemize}
  \item Sicherheit, Selbstschutz; Patientenanzahl; Umgebung, Ort,
Zeit
\item \EE{Trauma?}
\end{itemize}
\end{minipage}

&
\noindent\color{darkBlue}\itshape\begin{minipage}[t]{\linewidth}
\begin{itemize}
  \item Gefahrenzonen, Umstände,
weitere Resourcen erforderlich, Unfallmechanismus, Großschaden 
\end{itemize}
\end{minipage} 
\\\addlinespace\toprule\addlinespace
\noindent\TabHeading{Eindruck}
&
\begin{minipage}[t]{\linewidth}
\begin{itemize}
  \item \EE{Hautfarbe}, Gesichtsausdruck, \EE{Haltung},
  Sprache
  \item Sonstige \EE{Auffälligkeiten}
\end{itemize}
\end{minipage}
&
\noindent\color{darkBlue}\itshape\begin{minipage}[t]{\linewidth}
\begin{itemize}
  \item Blässe, Zyanose, 
  Schonhaltung, verwaschene Sprache,
  Krämpfe
  \item Trauma, insb. Hinweise auf WS-Verletzungen starke Blutungen
\end{itemize}
\end{minipage}
\\\addlinespace\midrule\addlinespace
\noindent\TabHeading{Bewusstsein}
&
\begin{minipage}[t]{\linewidth}
\begin{itemize}
  \item \EE{Bewusstseinsgrad}
\end{itemize}
\end{minipage}
&
\noindent\color{darkBlue}\itshape\begin{minipage}[t]{\linewidth}
\begin{itemize}
  \item Somnolenz, Sopor, Koma
\end{itemize}
\end{minipage}
\\\addlinespace\midrule\addlinespace
\noindent\TabHeading{Hauptbeschwerde}
&
\begin{minipage}[t]{\linewidth}
\begin{itemize}
  \item \EE{Berufungsgrund}, Schmerzen und andere dringliche
  \EE{Symptome}
\end{itemize}
\end{minipage}
&
\noindent\color{darkBlue}\itshape\begin{minipage}[t]{\linewidth}
\begin{itemize}
  \item Kritische Symptome
  \item Trauma
\end{itemize}
\end{minipage}
\\\addlinespace\toprule\addlinespace
\TabHeading{Atemwege}
&
\begin{minipage}[t]{\linewidth}
\begin{itemize}
  \item Inspektion (und Abtasten) von \EE{Mund}, \EE{Nase}
  und
  \EE{Hals}
  \item \EE{Atemgeräusche}
\end{itemize}
\end{minipage}
&
\noindent\color{darkBlue}\itshape\begin{minipage}[t]{\linewidth}
\begin{itemize}
  \item Pfeifendes, gurgelndes oder schnarchendes
  Atemgeräusch, Gestik
  \item[→] frei oder verlegt?
\end{itemize}
\end{minipage} 
\\\addlinespace\midrule\addlinespace
\TabHeading{Atmung}
&
\begin{minipage}[t]{\linewidth}
\begin{itemize}
  \item Schätzen: \EE{Atemfrequenz},\,\EE{-tiefe}
  \item Atemgeräusche
  \item Inspektion (und Abtasten) von \EE{Thorax} inkl.
  Atemprobe
  \item \E{Zeichen der
  Atemnot}: Hautfarbe, Bauch-/Thorax\-bewegungen,
  Anstrengung beim Atmen, Atemhilfsmuskulatur, \ldots
%   (\prettyref{topic:atemnot-zeichen-kinder})
  \item \Z{(Auskultation: Seitenvergleich, Lungenbasen)}
\end{itemize}
\end{minipage}
&
\noindent\color{darkBlue}\itshape\begin{minipage}[t]{\linewidth}
\begin{itemize}
  \item Zyanose
  \item Atemnot
  \item Tachy-/Bradypnoe
  \item Flache oder tiefe Atmung
  \item Atemmuster
  \item Fehlende, paradoxe, Schnappatmung 
  \item Heimsauerstoff, Atemhilfen
\end{itemize}
\end{minipage}
\\\addlinespace\midrule\addlinespace
\TabHeading{Kreislauf}
&
\begin{minipage}[t]{\linewidth}
\begin{itemize}
  \item \EE{Hautfarbe}, \EE{-temperatur} 
  \item Schweiß
  \item \EE{Rekap}-Zeit 
  \item \EE{Radialispuls}
  \item[{\tiny\color{Indigo}\ovalbox{Trauma}}] Inspektion
  und Abtasten von
  \begin{inparaenum}
    \item \EE{Bauch},
    \item \EE{Becken} und
    \item \EE{Oberschenkel}
  \end{inparaenum}
\end{itemize}
\end{minipage}
&
\noindent\color{darkBlue}\itshape\begin{minipage}[t]{\linewidth}
\begin{itemize}
  \item Sichtbare Blutungen
  \item Schockzeichen
  \item Rhythmus 
  
  (Tachy-/Bradykardie, Arrhythmie)
\end{itemize}
\end{minipage}
\\\addlinespace\midrule\addlinespace
\TabHeading{Weiters}
&
\begin{minipage}[t]{\linewidth}
\begin{itemize}
  \item \EE{Orientierung}
  \item \EE{Pupillen}
  \item \EE{Kraftprobe OE}
  \item[{\tiny\color{Indigo}\ovalbox{Trauma}}] Inspektion
  und Abtasten von
  \begin{inparaenum}
    \item \EE{Kopf} inkl. Ohren,
    \item \EE{Hals}
  \end{inparaenum}
  \item[{\tiny\color{Indigo}\ovalbox{Trauma}}] Kann der
  Patient \EE{Hände} und \EE{Füße} spüren und aktiv bewegen?
\end{itemize}
\end{minipage}
&
\noindent\color{darkBlue}\itshape\begin{minipage}[t]{\linewidth}
\begin{itemize}
  \item Entgleiste Vitawerte
  \item Verwirrtheit
  \item Pupillen: Abnorme \E{Weite}, \E{Differenz},
  \E{Entrundung}
  \item Motorische und sensible Ausfälle
  \item Kritische Verletzungen
\end{itemize}
\end{minipage}
\\\addlinespace\bottomrule
\end{tabularx}
\end{WidePar}

% \section{Sehen -- Hören -- Fühlen }
% \label{topic:sehen-hoeren-fuehlen}
% \TextSlide{Unsere eigenen Sinne}{
% Die Seh-, Hör- und Tastsinne sind bei der körperlichen
% Untersuchung unverzichtbar. Zuerst wird der Patient und sein
% Körper betrachtet. Wir versuchen, uns bereits bekannte Symptome wahrzunehmen.
% Hierbei können uns viele augenscheinliche Sachen auffallen: zum Beispiel
% \E{Fettleibigkeit}, \E{Zyanose}, \E{Blässe},
% \E{Schweiß}, \E{Fremdkörper}, \E{Verletzungen}, \E{Fehlstellungen},
% \ldots
% 
% Der nächste Schritt betrifft das Hören: Hier sind zum
% Beispiel die Atemgeräusche einzuordnen, die man eventuell schon ohne
% Hilfsmittel vernehmen kann (brodelndes Atemgeräusch beim
% Lungenödem, pfeifendes Atemgeräusch beim Asthmaanfall,
% röcheln bei Atemwegsverlegung, \ldots)
% 
% Wenn wir gesehen und gehört haben, können wir versuchen zu
% fühlen: Wir können den (Radialis-) Puls fühlen, können
% (stehende) Hautfalten beurteilen oder das Abdomen abtasten. Ebenso können wir
% die ungefähre Temperatur fühlen. }{
% \begin{itemize}
%   \item Zyanose
%   \item Blässe
%   \item Schweiß
%   \item Körperhaltung
%   \item Fremdkörper
%   \item Verletzungen, Fehlstellungen
%   \item Atemgeräusche
%   \item Puls
%   \item Stehende Hautfalten
%   \item Temperatur
%   \item Durchblutung des Nagelbetts
%   \item \ldots
% \end{itemize}
% }{}{}{}

\ChapterEnd